% 復習1　AIのハルシネーションを確かめる（記入例）
\session{1}{1}{40分}{AIのハルシネーションを確かめる}{出力の検証}{}

\goals{
  \item この授業でのAI利用の基本（大学の Copilot を使う、出力を検証する、使ったことを記録する）を説明できる
  \item AIがもっともらしい誤り（ハルシネーション）を出す場面を、自分で確かめられる
  \item 出力を使う前に何を確かめるかを、自分の言葉で挙げられる
}

\work{確認}{AI利用ガイドの「基本の考え方」（自分の言葉で4つ）}
\alines{4}{① 何を知りたいか（問い）は自分で決める。AIは視点を広げる相手。\par
② AIの出力は素材。そのまま貼らず、確かめて自分の言葉で書き直す。\par
③ 数値・固有名詞・文献は、統計や論文などの一次資料で確かめる。\par
④ どこでAIを使ったかを記録し、提出物に明記する。}

\work[12mm]{準備}{Copilot の動作確認}
\begin{achecks}
  \item 大学のアカウントで Copilot にサインインできた
  \item 簡単な質問をして、回答が返ってくることを確かめた
\end{achecks}

\work[70mm]{ワーク1}{AIに文献と数値を挙げさせ、確かめる}
\begin{promptbox}[ハルシネーションを体験する]
日本の中小企業における生成AIの導入状況について、参考文献を3つと具体的な数値を挙げて、300字で説明してください。
\end{promptbox}
\instr{AIが挙げた文献と数値が実在するか、CiNii Research、図書館の蔵書検索、官公庁の公式サイトで確かめる。}
\begin{wstabh}{colspec={X[3,l,m] Q[c,wd=14mm,m] X[2,l,m]},row{2-Z}={ht=10mm}}
  AIが挙げた文献・数値 & 実在 ○× & 確かめた場所・わかったこと \\
  \af{中小企業の生成AI活用を論じたとされる論文（著者・学会誌名・年つき）} & \ans{×} & \af{CiNii・J-STAGE とも該当なし。著者名でも見つからない} \\
  \af{中小企業庁の白書（年版つき）} & \ans{○} & \af{白書は実在。ただしAIが示した章の題名は、その年版の目次にない} \\
  \af{「中小企業の生成AI導入率は○\%」という数値} & \ans{×} & \af{AIは白書を出典としたが、白書の中に同じ数値が見つからない} \\
  \af{総務省の情報通信白書} & \ans{○} & \af{実在。生成AIの利用状況を扱う章がある。数値はAIの説明と違った} \\
\end{wstabh}

\work[50mm]{ワーク2}{自分がよく知っている話題で正誤を判定する}
\instr{アルバイト先の業界、出身地、部活、趣味など、自分が詳しく知っている話題について聞き、内容の正誤を判定する。}
\begin{wstabh}{colspec={X[2,l,m] X[3,l,m] Q[c,wd=14mm,m] X[2,l,m]},row{2-Z}={ht=12mm}}
  尋ねたこと & AIの答え（要点） & 正誤 ○△× & どこが違ったか \\
  \af{アルバイト先（金属加工）の業界の仕事の流れ} & \af{受注から設計、加工、検査、出荷までの流れ} & \ans{○} & \af{おおむね正しい。実際は検査が工程ごとにある} \\
  \af{出身地の地場産業の歴史} & \af{起源の時代と、代表的な企業名を断定的に挙げた} & \ans{△} & \af{起源の時代は市の資料と合うが、企業名の1つは存在しない} \\
  \af{地元の商店街の店舗数} & \af{具体的な数を挙げた} & \ans{×} & \af{商店街の公式サイトの数と大きく違う。出典も示されなかった} \\
\end{wstabh}

\work{まとめ}{気づいたこと・「出力を使う前のチェック」}
\instr{AI利用ガイドの「出力を使う前のチェック」も見ながら、自分の言葉で整理する。}
\alines{5}{AIは、実在する白書の名前に、存在しない数値や章を組み合わせることがある。半分正しいので見抜きにくい。\par
・文献は、題名だけでなく年版・章・ページまで確かめる\par
・数値は、AIが示した出典の中に同じ値があるかを確かめる\par
・よく知っている話題で一度試し、どの程度間違えるかの感覚を持つ\par
・自分で説明できない内容は使わない}

\begin{nexttime}
  \item 関心のあるテーマ（経営・情報・企業・働き方など）を3つ書き出す（復習2で使う）。「なぜ気になるのか」を一言添える。
\end{nexttime}
\begin{wstabh}{colspec={Q[c,wd=8mm,m] X[2,l,m] X[3,l,m]},row{2-Z}={ht=11mm}}
  & 関心のあるテーマ & なぜ気になるのか \\
  1 & \af{中小企業と生成AI} & \af{アルバイト先の社長が「AIを使ってみたいが、何に使えばよいかわからない」と話していたから} \\
  2 & \af{飲食店のセルフオーダー端末} & \af{よく行く店で導入され、店員の仕事の内容が変わったように見えるから} \\
  3 & \af{副業を認める企業の増加} & \af{就職先を考えるときに、働き方の選択肢として気になるから} \\
\end{wstabh}
\aimemoA{生成AIの導入状況の説明と、文献・数値の提示（ハルシネーションの確認）}
{「日本の中小企業における生成AIの導入状況について、参考文献を3つと具体的な数値を挙げて300字で」}
{出力は確認の練習の材料としてだけ使い、内容は採用しなかった。実在した白書2点は、復習3で自分で読む候補にした。}
{CiNii Research、J-STAGE、中小企業庁と総務省の公式サイト、商店街と市の公式サイト}

\begin{point}
  \item ハルシネーションは「全部うそ」より「実在する資料名＋存在しない数値」のような、半分正しい形で出ることが多い。資料が実在しても、中身まで確かめる。
  \item AIの出力は使うたびに変わる。実在しない文献が出なかった場合も、その結果を記録しておけばよい。「今回は出なかった」ことも検証の成果。
  \item 自分がよく知っている話題で試すと、AIの間違い方（固有名詞、数値、年）の傾向がつかめる。
  \item 関心テーマは、「なぜ気になるのか」まで書いておくと、復習2で問いに育てやすい。
\end{point}
